\documentclass[11pt]{article}
\usepackage[margin=1in]{geometry}
\usepackage{amsmath,amssymb,amsthm}
\usepackage{hyperref}
\usepackage{graphicx}

\newtheorem{theorem}{Theorem}
\newtheorem{lemma}{Lemma}
\newtheorem{remark}{Remark}

\newcommand{\R}{\mathbb R}
\newcommand{\N}{\mathbb N}
\newcommand{\HH}{\mathcal H}
\newcommand{\dist}{\operatorname{dist}}

\title{Endpoint unboundedness for the Riesz-potential operator in arXiv:2001.09677}
\author{Agent lane 19}
\date{2026-07-18}

\begin{document}
\maketitle

\section*{Source Question}

Hernandez--Semenov--Tradacete~\cite{HST} consider, for
\(0<\alpha,\lambda<1\), a closed Ahlfors \(\alpha\)-regular set
\(\Omega_\alpha\subset(0,1)\), equipped with
\(\mu=\mathcal H^\alpha|_{\Omega_\alpha}\), and the Riesz-potential operator
\[
  T_\lambda f(t)=\int_0^1 |t-u|^{-\lambda}f(u)\,du,
  \qquad t\in\Omega_\alpha .
\]
Their Theorem 10 proves boundedness and strict singular non-compactness on the
open line segment
\[
  \frac1q=\frac1\alpha\left(\frac1p-1+\lambda\right),
  \qquad
  1-\lambda < \frac1p < \frac{1-\lambda}{1-\alpha}.
\]
Remark 11 asks whether the endpoint operator is bounded and, if it is bounded,
whether it is strictly singular.

\begin{center}
\includegraphics[width=.92\linewidth]{figures/open_problem_crop.png}
\end{center}

\begin{remark}
The printed Remark 11 uses an undeclared symbol \(s\) in the first exponent.
The endpoint of Theorem 10, and the following sentence of the same remark,
identify the coherent Banach-space endpoint as
\[
  p=q=\frac{1-\alpha}{1-\lambda}.
\]
This is a Banach exponent precisely when \(\alpha<\lambda\), and this is the
case treated below.
\end{remark}

\section*{Theorem}

\begin{theorem}
Let \(0<\alpha<\lambda<1\), let
\(\Omega_\alpha\subset(0,1)\) be closed and Ahlfors \(\alpha\)-regular, and let
\(\mu=\mathcal H^\alpha|_{\Omega_\alpha}\). Put
\[
  p=\frac{1-\alpha}{1-\lambda}.
\]
Then \(T_\lambda\) is not bounded from \(L_p(0,1)\) to
\(L_p(\Omega_\alpha,\mu)\). Consequently, the endpoint asked about in
Remark 11 is outside \(L(T_\lambda)\), so the strict-singularity follow-up is
vacuous.
\end{theorem}

\section*{Proof Intuition}

At the endpoint, the adjoint potential generated by an Ahlfors-regular piece of
\(\Omega_\alpha\) has exactly the critical weak \(L^{p'}\) distribution near the
fractal set. Ahlfors regularity gives two competing powers: the potential grows
like \(r^{\alpha-\lambda}\), while the Lebesgue measure of the \(r\)-neighborhood
of the set is at least a constant times \(r^{1-\alpha}\). Since
\[
  p'=\frac{1-\alpha}{\lambda-\alpha},
\]
these powers cancel in the strong \(L^{p'}\) integral, leaving a logarithmic
divergence. Thus the adjoint cannot be bounded.

\section*{Proof}

Let \(p'=\frac{p}{p-1}\). Since
\[
  p=\frac{1-\alpha}{1-\lambda},
  \qquad
  p-1=\frac{\lambda-\alpha}{1-\lambda},
\]
we have
\[
  p'=\frac{1-\alpha}{\lambda-\alpha}.
\]

Assume for contradiction that
\[
  T_\lambda:L_p(0,1)\longrightarrow L_p(\Omega_\alpha,\mu)
\]
is bounded. Then its Banach-space adjoint is bounded
\[
  T_\lambda^*:L_{p'}(\Omega_\alpha,\mu)\longrightarrow L_{p'}(0,1),
\]
and for \(g\in L_{p'}(\Omega_\alpha,\mu)\),
\[
  T_\lambda^*g(x)=\int_{\Omega_\alpha}|x-t|^{-\lambda}g(t)\,d\mu(t)
  \quad\text{for a.e. }x\in(0,1).
\]
For nonnegative simple \(g\), this formula follows by applying Fubini to
truncated kernels and then using monotone convergence; boundedness would force
the displayed potential to belong to \(L_{p'}(0,1)\).

Choose \(t_0\in\Omega_\alpha\) and \(R>0\) so small that
\((t_0-R,t_0+R)\subset(0,1)\). Set
\[
  E=\Omega_\alpha\cap (t_0-R/2,t_0+R/2),
  \qquad
  E_0=\Omega_\alpha\cap (t_0-R/4,t_0+R/4),
  \qquad g=\mathbf 1_E .
\]
Then \(g\in L_{p'}(\Omega_\alpha,\mu)\). We will show that
\[
  P(x):=T_\lambda^*g(x)=\int_E |x-t|^{-\lambda}\,d\mu(t)
\]
does not belong to \(L_{p'}(0,1)\).

Let \(0<r<R/16\), and consider the \(r\)-neighborhood
\[
  N_r=\{x\in(0,1):\dist(x,E_0)<r\}.
\]
If \(x\in N_r\), choose \(y\in E_0\) with \(|x-y|<r\). For every
\(t\in\Omega_\alpha\cap(y-r,y+r)\) we have \(|x-t|<2r\). Since
\(\Omega_\alpha\cap(y-r,y+r)\subset E\), \(\Omega_\alpha\) is Ahlfors
\(\alpha\)-regular, and \(r<R/16\),
\[
  \mu(\Omega_\alpha\cap(y-r,y+r))\ge c r^\alpha,
\]
where \(c>0\) depends only on the regularity constants. Hence
\[
  P(x)\ge (2r)^{-\lambda}\mu(\Omega_\alpha\cap(y-r,y+r))
  \ge c_1 r^{\alpha-\lambda}
  \qquad (x\in N_r).
\]

We also need a lower bound for the Lebesgue measure of \(N_r\). Take a maximal
\(2r\)-separated subset \(\{y_1,\dots,y_m\}\) of
\(E_0\). The intervals
\((y_i-r,y_i+r)\) are disjoint and contained in \(N_r\), while maximality and
the upper Ahlfors estimate give
\[
  \mu(E_0)
  \le \sum_{i=1}^m \mu(\Omega_\alpha\cap(y_i-2r,y_i+2r))
  \le C m r^\alpha .
\]
The left-hand side is a positive constant depending on \(R\). Therefore
\(m\ge c_R r^{-\alpha}\), and
\[
  |N_r|\ge \sum_{i=1}^m |(y_i-r,y_i+r)|
  \ge c_2 r^{1-\alpha}.
\]

Write \(a=\lambda-\alpha>0\). The previous two estimates imply the distribution
lower bound
\[
  \big|\{x\in(0,1):P(x)\ge c_1 r^{-a}\}\big|
  \ge c_2 r^{1-\alpha}
  \qquad (0<r<R/16).
\]
Since \(p'=(1-\alpha)/a\), this is exactly the critical weak \(L^{p'}\) decay.
To see that the strong \(L^{p'}\) norm diverges, take \(r_n=2^{-n}r_0\), with
\(r_0<R/16\), and \(\tau_n=c_1 r_n^{-a}\). By the layer-cake formula,
\[
  \int_0^1 P(x)^{p'}\,dx
  =
  p'\int_0^\infty \tau^{p'-1}
  |\{x:P(x)>\tau\}|\,d\tau .
\]
For \(\tau\in[\tau_n,\tau_{n+1}]\),
\[
  |\{P>\tau\}|\ge |\{P>\tau_{n+1}\}|
  \ge c\,r_{n+1}^{1-\alpha}
  = c'\tau_{n+1}^{-p'}.
\]
Thus each interval contributes at least
\[
  c'p'\tau_{n+1}^{-p'}\int_{\tau_n}^{\tau_{n+1}}\tau^{p'-1}\,d\tau
  =
  c''\left(1-\left(\frac{\tau_n}{\tau_{n+1}}\right)^{p'}\right),
\]
which is bounded below by a positive constant independent of \(n\). Summing over
\(n\) gives
\[
  \int_0^1 P(x)^{p'}\,dx=\infty.
\]
Therefore \(T_\lambda^*g\notin L_{p'}(0,1)\), contradicting boundedness of the
adjoint. Hence \(T_\lambda\) is not bounded at the endpoint.

\section*{Verification Notes}

The proof uses only the two-sided Ahlfors regularity of \(\Omega_\alpha\) and
elementary duality. It does not rely on any endpoint trace theorem. The condition
\(\alpha<\lambda\) is exactly what makes the endpoint exponent
\((1-\alpha)/(1-\lambda)\) larger than \(1\); if \(\lambda=\alpha\) the endpoint
is \(p=1\), and if \(\lambda<\alpha\) it is below the Banach range.

The source already proves noncompactness for the open segment. The present
argument shows that the missing endpoint is not merely noncompact; it is not
bounded at all.

\section*{Novelty Check}

Local index searches for \texttt{2001.09677}, \texttt{strictly singular non compact},
\texttt{Riesz potential}, \texttt{endpoint}, and related \(L_p\)-characteristic terms found no
prior packet for this source question. Web searches for exact endpoint phrases,
the source title together with \texttt{T\_lambda}, and Adams/Riesz trace-inequality
keywords found broad trace-inequality literature, including classical endpoint
caveats for \(p=q\), but no direct answer to Remark 11 of
Hernandez--Semenov--Tradacete.

\section*{Human Review Recommendation}

Recommended for review as a likely valid full answer to the endpoint boundedness
question in Remark 11, interpreting the remark's undeclared \(s\) as the
Ahlfors exponent \(\alpha\) dictated by Theorem 10.

\begin{thebibliography}{9}
\bibitem{HST}
F. L. Hernandez, E. M. Semenov, and P. Tradacete,
\emph{Strictly singular non-compact operators between \(L_p\) spaces},
arXiv:2001.09677.
\end{thebibliography}

\end{document}
